% !TeX root = booklet-02.tex
\chapter*{مرور پیش‌نیازهای لازم}
\addcontentsline{toc}{chapter}{مرور پیش‌نیازهای لازم}
\section*{ولتاژ، میدان و جهت حرکت}
ولتاژ، اختلاف پتانسیل دو نقطه و انرژی به ازای واحد بار است. در نماد
\(V_{DS}=V_D-V_S\)، ترتیب دو حرف، ترتیب تفریق را تعیین می‌کند. اگر
\(V_G=\qty{1.2}{\volt}\)، \(V_S=\qty{0.2}{\volt}\) و
\(V_D=\qty{0.5}{\volt}\) باشند، داریم \(V_{GS}=\qty{1.0}{\volt}\) و
\(V_{DS}=\qty{0.3}{\volt}\). شرط هدایت با اختلاف پتانسیل \lr{Gate} و
\lr{Source} سنجیده می‌شود \cite[بخش \secnum{۲.۱}]{weste}.

بزرگی متوسط میدان طولی در یک فاصله \(L\)، تقریباً \(V_{DS}/L\) است.
برای مثال، ولتاژ \(\qty{1}{\volt}\) روی طول \(\qty{50}{\nano\metre}\)
میدان متوسط \(\qty{2e7}{\volt\per\metre}\) ایجاد می‌کند. در امتداد کانال،
پتانسیل و میدان می‌توانند با مکان تغییر کنند. در \lr{nMOS} با \(V_D>V_S\)،
الکترون‌ها از \lr{Source} به \lr{Drain} حرکت می‌کنند و جریان قراردادی جهت
مخالف دارد. در روابط سرعت این جزوه، \(v\) و \(E\) بزرگی‌های مثبت‌اند
\cite[بخش‌های \secnum{۲.۱} و \secnum{۲.۴.۱}]{weste}.

\section*{از بار به جریان و از هندسه به ضریب جریان}
جریان، مقدار بار عبوری در واحد زمان است. اگر در هر متر کانال، بزرگی بار
متحرک \(q'\) وجود داشته باشد و حامل‌ها با سرعت متوسط \(v\) جابه‌جا شوند،
در یک ثانیه بار موجود در طول \(v\) متر از یک مقطع عبور می‌کند. بنابراین
رابطه جریان و کنترل واحد آن چنین‌اند.
\[
 I=q'v,\qquad
 \mathrm{C\,m^{-1}}\times\mathrm{m\,s^{-1}}=\mathrm{A}.
\]
در مدل بار کانال، رابطه زیر به مکان وابسته است.
\[
 q'(x)=WC_{\mathrm{ox}}[V_{GT}-V(x)].
\]
\cite[بخش \secnum{۲.۲}]{weste}. نماد پریم در \(q'\) به معنای «بر واحد طول»
است و برای تشخیص واحد این کمیت به کار می‌رود.

ظرفیت اکسید بر واحد سطح از نسبت گذردهی به ضخامت به دست می‌آید.
\[
 C_{\mathrm{ox}}=\frac{\varepsilon_{\mathrm{ox}}}{t_{\mathrm{ox}}}.
\]
واحد این کمیت \(\unit{\farad\per\metre\squared}\) است. عرض و طول کانال به‌ترتیب
\(W,L\) و ضخامت اکسید \(t_{\mathrm{ox}}\) هستند. با \lr{Mobility} ثابت،
ضریب جریان با واحد \(\unit{\ampere\per\volt\squared}\) تعریف می‌شود.
\[
 \beta=\mu C_{\mathrm{ox}}\frac{W}{L}.
\]
در مدل ایده‌آل و برای \(V_{GT}>0\)، جریان و شرط هر ناحیه چنین‌اند.
\[
 \begin{aligned}
 I_{DS}&=\beta(V_{GT}V_{DS}-V_{DS}^2/2),
 &0\le V_{DS}&<V_{GT},\\
 I_{DS}&=\beta V_{GT}^2/2,
 &V_{DS}&\ge V_{GT}.
 \end{aligned}
\]
سطر اول ناحیه \lr{Linear} و سطر دوم ناحیه \lr{Saturation} را بیان می‌کند.
مرز این دو ناحیه \(V_{DS}=V_{GT}\) است
\cite[رابطه \secnum{۲.۱۰}]{weste}. با تغییر مدل، شرط مرز نیز باید از همان
مدل به دست آید.

\section*{تبدیل واحدها و کار با توان}
روابط عددی فقط با واحدهای سازگار نتیجه درست می‌دهند. تبدیل‌های مورد نیاز
در این گام چنین‌اند
\[
 \qty{1}{\centi\metre}=\qty{e-2}{\metre},\qquad
 \qty{1}{\centi\metre\squared}=\qty{e-4}{\metre\squared},\qquad
 \qty{1}{\nano\metre}=\qty{e-9}{\metre}.
\]
در نتیجه، تبدیل \lr{Mobility} و سرعت می‌دهد
\[
 \begin{aligned}
 \qty{96}{\centi\metre\squared\per\volt\per\second}
 &=\qty{0.0096}{\metre\squared\per\volt\per\second},\\
 \qty{e7}{\centi\metre\per\second}&=\qty{e5}{\metre\per\second}.
 \end{aligned}
\]
عامل تبدیل طول هنگام تبدیل مساحت به توان دو می‌رسد.

توان غیرصحیح نیز مانند توان صحیح، روی یک عدد بی‌بعد اعمال می‌شود.
برای محاسبه \(a^{1.85}\) با \(a>0\)، می‌توان از
\(a^{1.85}=\exp(1.85\ln a)\) استفاده کرد. \(\ln\) لگاریتم طبیعی و
\(\exp\) تابع نمایی با پایه \(e\) است. برای نمونه،
\(2^{1.85}\approx3.61\). در استخراج توان تجربی، از این خاصیت استفاده
می‌کنیم که \(\ln(a^b)=b\ln a\) و \(\ln(a/b)=\ln a-\ln b\).

\section*{نسبت، شیب و میانگین پیوسته}
در کسر \(I/I_0=1/(1+x)\)، اگر \(x=0.25\) باشد، نسبت \(0.8\) است.
پس مقدار جدید ۸۰ درصد مقدار مرجع و کاهش آن ۲۰ درصد است. مقدار مثبت
\(x\) در مخرج نسبت را کوچک‌تر از یک می‌کند. برای \(|x|\ll1\)، تقریب
\(1/(1+x)\approx1-x\) با حذف جمله مرتبه دوم به دست می‌آید. هویت دقیق
\(1/(1+x)=1-x+x^2/(1+x)\) اندازه جمله حذف‌شده را مشخص می‌کند.

شیب نمودار جریان برحسب ولتاژ برابر \(\Delta I/\Delta V\) و واحد آن
\(\unit{\ampere\per\volt}=\unit{\siemens}\) است. برای خط
\(I=b+mV\)، \(m\) شیب و \(b\) عرض از مبدأ است. اگر تغییر جریان در بازه
ولتاژ یکنواخت نباشد، مشتق \(\mathrm{d}I/\mathrm{d}V\) شیب محلی را بیان
می‌کند. برای تابع خطی، این شیب در همه نقاط برابر \(m\) است.

در استخراج مدل جریان، انتگرال را برای جمع سهم‌های کوچک کانال به کار
می‌بریم. وقتی \(I\) در تمام مقاطع ثابت باشد،
\(\int_0^L I\,\mathrm{d}x=IL\). همچنین
\(\int_0^D(A-V)\,\mathrm{d}V=AD-D^2/2\)، زیرا انتگرال جمله ثابت \(A\)
برابر \(AV\) و انتگرال \(V\) برابر \(V^2/2\) است. ارزیابی در دو کران
یعنی مقدار عبارت در \(V=D\) منهای مقدار آن در \(V=0\).
